\documentclass[aspectratio=169]{beamer}
\usepackage{graphicx}
\usepackage{booktabs}
\usepackage{amsmath}

\usetheme{Madrid}
\usecolortheme{default}

\title{Congressional Trade Dataset Preparation}
\subtitle{Methodology for Data Cleaning, Enrichment, and Labeling}
\author{Eugene Sy}
\institute{National Taiwan University}
\date{\today}

\begin{document}

\begin{frame}
    \titlepage
\end{frame}

\begin{frame}
    \frametitle{Project Objective & Overview}
    \begin{columns}
        \begin{column}{0.6\textwidth}
            \textbf{Objective:} To construct a clean, labeled dataset from raw congressional financial disclosures to enable analysis of trading behavior and performance.

            \vspace{1em}
            \textbf{Workflow:}
            \begin{enumerate}
                \item Ingest raw transaction data.
                \item Enrich data with market benchmarks (S&P 500) and historical prices.
                \item Engineer features, including various return metrics.
                \item \textbf{Clean and standardize} the dataset to ensure quality.
                \item \textbf{Construct a binary label} to classify trade performance.
            \end{enumerate}
        \end{column}
        \begin{column}{0.4\textwidth}
            \textbf{Initial Data:}

            \texttt{v5_transactions.csv}
            \begin{itemize}
                \item Rows: 35,365
                \item Columns: 26
            \end{itemize}
            \vspace{1em}
            \textbf{Final Data:}

            \texttt{v9_transactions.csv}
            \begin{itemize}
                \item Rows: 32,416
                \item Columns: 47
            \end{itemize}
        \end{column}
    \end{columns}
\end{frame}

\begin{frame}
    \frametitle{Core Task 1.1: Data Cleaning - Structural & Missing Data}
    A significant portion of the project involved cleaning inconsistent and missing data. A total of \textbf{2,949 rows} were removed during this process.

    \textbf{Key Cleaning Steps (Part 1):}
    \begin{itemize}
        \item \textbf{Duplicate Removal:} Checked for and removed 0 duplicate rows, ensuring unique transaction entries.
        \item \textbf{Column Pruning:} Removed columns with high nullity or low analytical value:
        \begin{itemize}
            \item \texttt{Short} (35,314 NaNs), \texttt{Description} (32,785 NaNs), \texttt{Subholding} (14,042 NaNs), \texttt{District}.
        \end{itemize}
        \item \textbf{Critical Missing Data Handling:}
        \begin{itemize}
            \item Dropped rows where the critical \texttt{Label_1M} was \texttt{NaN} (e.g., due to insufficient historical data).
            \item Filled categorical NaNs in \texttt{Committees}, \texttt{Industry}, and \texttt{Sector} with `Unknown'.
        \end{itemize}
    \end{itemize}
\end{frame}

\begin{frame}
    \frametitle{Core Task 1.2: Data Cleaning - Standardization}
    Continuing the data cleaning process, focusing on standardizing categorical and structural features.

    \textbf{Key Cleaning Steps (Part 2):}
    \begin{itemize}
        \item \textbf{Transaction Standardization:}
        \begin{itemize}
            \item Consolidated `Sale (Full)' and `Sale (Partial)' into a single `Sale' category.
            \item Removed 266 `Exchange' transactions, as they do not represent typical buy/sell signals.
        \end{itemize}
        \item \textbf{Party Affiliation Standardization:} Mapped abbreviations (D, R, I) to full names (Democrat, Republican, Independent).
        \item \textbf{Trade Size Standardization:} Converted raw `Trade_Size_USD` values into standardized categorical ranges (e.g., '1 - 1000', '1001 - 15000').
        \item \textbf{Unique Transaction IDs:} Assigned a unique `ID` to each transaction for easier tracking and referencing.
    \end{itemize}
\end{frame}

\begin{frame}
    \frametitle{Missing Value Analysis}
    Understanding and addressing missing values was crucial for data quality.

    \textbf{Initial Missing Values (Top 5):}
    \begin{itemize}
        \item \texttt{Short}: 35,314
        \item \texttt{Comments}: 34,525
        \item \texttt{Description}: 32,785
        \item \texttt{Subholding}: 14,042
        \item \texttt{District}: 8,227
    \end{itemize}
    \vspace{1em}
    \textbf{Final Missing Values (Post-Cleaning):}
    \begin{itemize}
        \item \texttt{Comments}: 31,660 (Retained as non-critical for core analysis)
        \item \texttt{SPY_6Months}, \texttt{Label_6M}, \texttt{Excess_Return_6M}, \texttt{SPY_Return_6M}: \textasciitilde1,539 (Due to future date projections or data availability)
        \item \texttt{SPY_3Months}, \texttt{Label_3M}, \texttt{Excess_Return_3M}, \texttt{SPY_Return_3M}: \textasciitilde748 (Similar reasons as 6-month data)
    \end{itemize}
    \vspace{1em}
    \textit{The reduction in critical missing values ensures a robust dataset for modeling.}
\end{frame}

\begin{frame}
    \frametitle{Visualizing Cleaning: Transaction & Party Standardization}
    \begin{columns}
        \begin{column}{0.5\textwidth}
            \textbf{Transaction Consolidation}

            \includegraphics[width=\textwidth]{../data/comparison_results/trade_size_distribution.png}
        \end{column}
        \begin{column}{0.5\textwidth}
            \textbf{Party Affiliation Mapping}

            \includegraphics[width=\textwidth]{../data/comparison_results/party_distribution.png}
        \end{column}
    \end{columns}
\end{frame}

\begin{frame}
    \frametitle{Core Task 2: Building the Trading Label}
    The primary goal is to create a binary label representing a ``successful'' trade. This requires defining success based on market outperformance.

    \vspace{1em}
    \textbf{Step 1: Calculate Excess Return}
    \begin{itemize}
        \item First, we calculate the stock's return and the benchmark (SPY) return over a given period (1, 3, 6 months).
        \item Then, we define Excess Return as the difference:
        \[
        \text{Excess Return} = R_{\text{stock}} - R_{\text{SPY}}
        \]
    \end{itemize}
    \vspace{1em}
    \textbf{Step 2: Define the Label Logic}
    \begin{itemize}
        \item We hypothesize a ``good'' trade outperforms the market by a significant margin, set here at \textbf{6\%}.
        \item For a \textbf{Purchase}: Label = 1 if \(\text{Excess Return} > 0.06
        \item For a \textbf{Sale}: Label = 1 if \(\text{Excess Return} \leq 0.06
    \end{itemize}
\end{frame}

\begin{frame}
    \frametitle{Visualizing the Results: Returns and Labels}
    \begin{columns}
        \begin{column}{0.5\textwidth}
            \textbf{1-Month Excess Return Distribution}
            \begin{itemize}
                \item The distribution is centered near zero, which is expected.
                \item It has long tails, indicating trades with significant outperformance or underperformance.
            \end{itemize}

            \includegraphics[width=\textwidth]{../data/comparison_results/excess_return_1m_distribution.png}
        \end{column}
        \begin{column}{0.5\textwidth}
            \textbf{1-Month Label Distribution}
            \begin{itemize}
                \item The resulting labels are almost perfectly balanced.
                \item This is ideal for training an unbiased machine learning model.
            \end{itemize}

            \includegraphics[width=\textwidth]{../data/comparison_results/label_1m_distribution.png}
        \end{column}
    \end{columns}
\end{frame}

\begin{frame}
    \frametitle{Conclusion & Next Steps}
    \textbf{Accomplishments:}
    \begin{itemize}
        \item A robust, automated pipeline for processing and cleaning financial disclosure data has been established.
        \item A clear, hypothesis-driven trading label (Label_1M) has been engineered, resulting in a well-balanced dataset.
        \item The final dataset is significantly enriched and standardized, making it suitable for advanced analysis.
    \end{itemize}
    \vspace{1em}
    \textbf{Next Steps:}
    \begin{itemize}
        \item Utilize this labeled dataset to train classification models (e.g., Logistic Regression, Gradient Boosting) to predict successful trades.
        \item Explore feature importance to understand what factors (e.g., party, committee assignments, sector) are most predictive of trade performance.
        \item Analyze the model's predictions to identify patterns in congressional trading.
    \end{itemize}
\end{frame}

\end{document}